\chapter{A Capstone: the Networked Logger}
\label{ch:capstone}\index{capstone}

Every chapter so far has built one subsystem in isolation: a sensor driver, a
filesystem, a TLS client, a periodic task. A real product is what you get when those
pieces compose. This closing chapter of Part~IV sketches one such product, a
\emph{networked environmental logger}, not as a new driver but as an assembly of the
ones the book already ships. It reads a sensor on a schedule, timestamps and stores
each reading, and uploads batches over HTTPS, and it does all of that inside the
static, bounded, type-checked envelope the rest of the book has been building toward.

\section{The shape of the product}

The logger has four jobs, and each is a subsystem you have already met:

\begin{description}[style=nextline, leftmargin=0pt]
\item[Measure] read temperature and humidity from an SHT41 over I2C
  (Appendix~\ref{app:external}), on a fixed period.
\item[Timestamp] keep wall-clock time, set once from an NTP server over the network
  (Chapter~\ref{ch:networking}) and carried across sleeps in the RTC.
\item[Store] append each reading to a file on an ext4 volume, on an SD card or the
  on-board SPI-NOR flash (Chapter~\ref{ch:storage}), so nothing is lost to a power
  cut.
\item[Upload] batch the stored readings and POST them to a server over TLS~1.3
  (Chapter~\ref{ch:tls}), then mark them sent.
\end{description}

Nothing here is new code in the sense that matters. Each job has a running example in
the tree; the capstone is the wiring.

\section{The task structure}\index{task!periodic}

The skeleton is a single library-level periodic task, exactly the shape of
Chapter~\ref{ch:tasking}. It owns the measure-store rhythm; the upload runs less
often, gated by a count. Because the task is declared at library level, it is legal
under the Jorvik profile, and because its period is computed from an absolute time it
holds its rate with no cumulative drift:

\begin{lstlisting}
task body Logger is
   Next : Time := Clock;
begin
   Bring_Up;                       --  mount the volume, open the NIC, set the clock
   loop
      delay until Next;
      declare
         R : constant Reading := Sensor.Read;          --  SHT41 over I2C
      begin
         Store.Append (Stamp => Now, Sample => R);     --  one line to ext4
         if Store.Pending >= Batch then
            Upload.Flush;                              --  HTTPS POST, then mark sent
         end if;
      end;
      Next := Next + Period;
   end loop;
end Logger;
\end{lstlisting}

The four collaborators (\texttt{Sensor}, \texttt{Store}, \texttt{Upload}, and the
clock) are packages with private state, each guarding its own resource the way
Chapter~\ref{ch:packages} described. The task never touches a register or a socket
directly; it speaks only to those contracts.

\section{Measure and timestamp}

\texttt{Sensor.Read} is a thin wrapper over \texttt{ESP32S3.SHT41}: a single I2C
transaction that returns a temperature and a humidity in fixed-point, with no float
in sight. Time is the only state that must outlive a reset. It is set once at start,
from NTP if the network is up, and then carried in the RTC so that a reset or a deep
sleep does not lose the wall clock (Chapter~\ref{ch:lowpower}). Each reading is
stamped from that clock as it is taken.

\section{Store: durable by construction}\index{ext4}

\texttt{Store.Append} formats one reading as a line and appends it to a file on the
mounted ext4 volume. The filesystem earns its keep here. The write-back cache makes
the per-reading cost small, and the journal makes a power cut at the wrong instant
survivable: on the next boot the volume replays cleanly rather than corrupting the
log (Chapter~\ref{ch:storage}). \texttt{Store.Pending} counts the readings written
but not yet uploaded, which is the only handshake the upload step needs.

\section{Upload: batched and authenticated}\index{TLS}

\texttt{Upload.Flush} opens a socket through the chip-neutral network stack, runs a
TLS~1.3 handshake with full certificate-chain validation against a pinned root, and
POSTs the pending readings as one request body (Chapter~\ref{ch:tls}). Batching is
deliberate: a handshake costs tens of milliseconds of curve arithmetic, so amortising
it over many readings is what keeps the average cost low. Only when the server has
acknowledged the batch does \texttt{Store} mark those readings sent, so a dropped
connection loses nothing; the next flush simply re-sends them.

\section{Resilience, for free}

Three of the book's later chapters turn the demonstrator into something you could
deploy, and each is a drop-in rather than a rewrite:

\begin{itemize}
\item \textbf{Failover.} With a second NIC, the routing table and interface pinning
  (Chapter~\ref{ch:networking}) let \texttt{Upload} ride whichever link is up,
  fail-closed if none is, with no change to the logger task.
\item \textbf{Field updates.} The OTA design (Chapter~\ref{ch:ota}) lets the same
  board take a signed firmware update over the network it already has open, with a
  safe rollback if the new image fails to confirm.
\item \textbf{Battery.} When readings are minutes apart, the logger sleeps between
  them (Chapter~\ref{ch:lowpower}): it stores a reading, records its phase in
  retained RTC memory, and deep-sleeps on a timer, waking to take the next.
\end{itemize}

\section{The whole budget, accounted for}\index{static memory}

The point of the capstone is not that it does a lot, but that it does it inside the
envelope the book has insisted on throughout. The task set is fixed, so every stack
is sized at link time. There is no hidden allocation in the steady state: the sensor
returns by value, the log line is a bounded buffer, the TLS session and the socket
storage are reserved once. The whole footprint, code and data and every stack, reads
off the two static-analysis commands of Chapter~\ref{ch:static-mem}, and the
worst-case stack depth of the logger task is a number you can compute and check
against its allotment. A networked, encrypted, power-loss-tolerant data logger, with
no operating system and no C library under it, that still fits in a budget you can
prove. That is the thesis of this book, built once into a single application.
